Unless otherwise stated, $(M,g)$ is a Riemannian manifold without boundary;
statements using the global Riemannian distance assume that $M$ is connected.

\section{Geodesics and Minimizing Curves}

An admissible curve $\gamma$ is \emph{minimizing} if
$L_g(\gamma)\leq L_g(\widetilde\gamma)$ for every admissible curve with the
same endpoints.  When $M$ is connected, this is equivalent to
$L_g(\gamma)=d_g(\gamma(a),\gamma(b))$.

\section{Families of Curves}

\begin{definition}[One-Parameter Family of Curves]
\label{def:one-parameter-family-of-curves}

For intervals $I,J\subseteq\mathbb R$, a \emph{one-parameter family of
curves} is a continuous map $\Gamma:J\times I\to M$.  Its main and transverse
curves are
\[
\Gamma_s(t)=\Gamma(s,t),
\qquad
\Gamma^{(t)}(s)=\Gamma(s,t),
\]
respectively.  When $\Gamma$ is continuously differentiable, write
\[
\partial_t\Gamma(s,t)=(\Gamma_s)'(t),
\qquad
\partial_s\Gamma(s,t)=(\Gamma^{(t)})'(s).
\]
A vector field along $\Gamma$ is a continuous map $V:J\times I\to TM$ with
$V(s,t)\in T_{\Gamma(s,t)}M$.
\end{definition}

An \emph{admissible family of curves} is a one-parameter family
$\Gamma:J\times[a,b]\to M$ such that $J$ is open, every main curve is
admissible, and some partition $(a_0,\ldots,a_k)$ of $[a,b]$ makes $\Gamma$
smooth on each $J\times[a_{i-1},a_i]$.  Such a partition is an
\emph{admissible partition} for the family.

\begin{definition}[Variations]
\label{def:variations-proper-variations}

Let $\gamma:[a,b]\to M$ be admissible.  A \emph{variation} of $\gamma$ is an
admissible family $\Gamma:J\times[a,b]\to M$, where $J$ is an open interval
containing $0$ and $\Gamma_0=\gamma$.  It is \emph{proper} if
\[
\Gamma_s(a)=\gamma(a),
\qquad
\Gamma_s(b)=\gamma(b)
\]
for every $s\in J$.
\end{definition}

A vector field along an admissible family is \emph{piecewise smooth} if it is
continuous and smooth on the rectangles of some admissible partition.
The field $\partial_s\Gamma$ is piecewise smooth, although
$\partial_t\Gamma$ may jump at partition points.  For a variation of
$\gamma$, its \emph{variation field} is
$V(t)=\partial_s\Gamma(0,t)$.  A field along $\gamma$ is \emph{proper} if
$V(a)=V(b)=0$; every proper variation has a proper variation field.

\begin{lemma}[Existence of a Variation with Prescribed Field]
\label{lem:variation-field-existence}
\dcref{ch6:6.1}

If $\gamma$ is admissible and $V$ is a piecewise smooth vector field along
$\gamma$, then some variation of $\gamma$ has variation field $V$.  If $V$ is
proper, the variation may be chosen proper.
\end{lemma}

\begin{proof}
\sketch
Set
\[
\Gamma(s,t)=\exp_{\gamma(t)}(sV(t)).
\]
Compactness of $[a,b]$ gives $\varepsilon>0$ for which this is defined on
$(-\varepsilon,\varepsilon)\times[a,b]$.  It is continuous globally and
smooth on each rectangle on which $V$ is smooth, and
$d(\exp_{\gamma(t)})_0=\operatorname{id}$ shows that its variation field is
$V$.  If $V(a)=V(b)=0$, both endpoint curves are constant, so the variation
is proper.
\end{proof}

For a piecewise smooth field $V$ along an admissible family, write $D_tV$ and
$D_sV$ for covariant differentiation along the main and transverse curves,
respectively, wherever the derivatives are defined.

\begin{lemma}[Symmetry Lemma]
\label{lem:symmetry-lemma}
\dcref{ch6:6.2}

Let $\Gamma:J\times[a,b]\to M$ be an admissible family.  On every rectangle
$J\times[a_{i-1},a_i]$ on which $\Gamma$ is smooth,
\[
D_s\partial_t\Gamma=D_t\partial_s\Gamma.
\]
\end{lemma}

\begin{proof}
\sketch
In local coordinates, writing $\Gamma=(x^1,\ldots,x^n)$ gives
\[
D_s\partial_t\Gamma
=\left(\frac{\partial^2x^k}{\partial s\partial t}
+\frac{\partial x^i}{\partial t}\frac{\partial x^j}{\partial s}
\Gamma^k_{ij}\right)\partial_k
\]
and
\[
D_t\partial_s\Gamma
=\left(\frac{\partial^2x^k}{\partial t\partial s}
+\frac{\partial x^i}{\partial s}\frac{\partial x^j}{\partial t}
\Gamma^k_{ij}\right)\partial_k.
\]
Equality follows from commutation of mixed partial derivatives and
$\Gamma^k_{ij}=\Gamma^k_{ji}$.
\end{proof}

\section{Minimizing Curves Are Geodesics}

\begin{theorem}[First Variation Formula]
\label{thm:first-variation-formula}
\uses{lem:symmetry-lemma}
\dcref{ch6:6.3}


Let $\gamma:[a,b]\to M$ be a unit-speed admissible curve, let
$\Gamma:J\times[a,b]\to M$ be a variation of $\gamma$, and let $V$ be its
variation field.  Then $s\mapsto L_g(\Gamma_s)$ is smooth and
\begin{equation}
\frac{d}{ds}\bigg|_{s=0} L_g(\Gamma_s)
=-\int_a^b\langle V,D_t\gamma'\rangle\,dt
-\sum_{i=1}^{k-1}\langle V(a_i),\Delta_i\gamma'\rangle
+\langle V(b),\gamma'(b)\rangle-\langle V(a),\gamma'(a)\rangle, \tag{6.1}
\end{equation}
where $(a_0,\ldots,a_k)$ is an admissible partition for $\gamma$ and, for
$i=1,\ldots,k-1$,
\[
\Delta_i\gamma'=\gamma'(a_i^+)-\gamma'(a_i^-).
\]
For a proper variation, this reduces to
\begin{equation}
\frac{d}{ds}\bigg|_{s=0} L_g(\Gamma_s)
=-\int_a^b\langle V,D_t\gamma'\rangle\,dt
-\sum_{i=1}^{k-1}\langle V(a_i),\Delta_i\gamma'\rangle. \tag{6.2}
\end{equation}
\end{theorem}

\begin{proof}
\sketch
On a rectangle of smoothness, set
$T=\partial_t\Gamma$ and $S=\partial_s\Gamma$.  Differentiation under the
integral and the symmetry lemma give
\begin{align*}
\frac{d}{ds}L_g\bigl(\Gamma_s|_{[a_{i-1},a_i]}\bigr)
&=\int_{a_{i-1}}^{a_i}\frac{1}{|T|}
\langle D_sT,T\rangle\,dt \\
&=\int_{a_{i-1}}^{a_i}\frac{1}{|T|}
\langle D_tS,T\rangle\,dt. \tag{6.3}
\end{align*}
At $s=0$, $T=\gamma'$ has unit length and $S=V$.  Metric compatibility and
integration by parts yield
\[
\int_{a_{i-1}}^{a_i}\langle D_tV,\gamma'\rangle\,dt
=\langle V(a_i),\gamma'(a_i^-)\rangle
-\langle V(a_{i-1}),\gamma'(a_{i-1}^+)\rangle
-\int_{a_{i-1}}^{a_i}\langle V,D_t\gamma'\rangle\,dt.
\]
Summing over the partition produces the jump and endpoint terms in (6.1);
properness removes the endpoint terms and gives (6.2).  Smoothness follows
from differentiating the finitely many smooth integrals over compact
subintervals.
\end{proof}

Every admissible curve admits a unit-speed parametrization, and length is
invariant under reparametrization.

\begin{theorem}[Minimizing Curves Are Geodesics]
\label{thm:minimizing-geodesic-unit-speed}
\uses{lem:symmetry-lemma}
\dcref{ch6:6.4}


Every minimizing curve in a Riemannian manifold is a geodesic after
unit-speed reparametrization.
\end{theorem}

\begin{proof}
\sketch
Let $\gamma:[a,b]\to M$ be minimizing and unit speed.  For every proper
variation, $s\mapsto L_g(\Gamma_s)$ has a minimum at $0$; hence (6.2)
vanishes for every proper piecewise smooth field $V$.

On a smooth partition interval choose a bump function $\varphi$ that is
positive in its interior and vanishes outside it.  Substitution of
$V=\varphi D_t\gamma'$ into (6.2) gives
\[
0=-\int\varphi|D_t\gamma'|^2\,dt,
\]
so $D_t\gamma'=0$ on every subinterval.  At a partition point $a_i$, choose
a proper piecewise smooth field satisfying
$V(a_i)=\Delta_i\gamma'$ and $V(a_j)=0$ for $j\ne i$.  Formula (6.2) then
gives $-|\Delta_i\gamma'|^2=0$.  Thus the velocity has no jumps, and uniqueness
of geodesics joins the smooth pieces into one geodesic.
\end{proof}

The proof of Theorem~\ref{thm:minimizing-geodesic-unit-speed} uses only
stationarity of the length functional.  A unit-speed admissible curve is a
\emph{critical point} of $L_g$ if the derivative of $L_g(\Gamma_s)$ at $s=0$
vanishes for every proper variation.  Thus
Theorem~\ref{thm:minimizing-geodesic-unit-speed} has the following
strengthening.

\begin{corollary}[Critical Curves]
\label{cor:critical-point-geodesic}
\uses{thm:minimizing-geodesic-unit-speed}
\dcref{ch6:6.5}


A unit-speed admissible curve is a critical point of $L_g$ if and only if it
is a geodesic.
\end{corollary}

\begin{proof}
\sketch
For a critical curve, the proof of
Theorem~\ref{thm:minimizing-geodesic-unit-speed} applies verbatim.  Conversely,
for a geodesic the integral in (6.2) vanishes because $D_t\gamma'=0$, and the
jump sum vanishes because $\gamma'$ is smooth.
\end{proof}

\section{Geodesics Are Locally Minimizing}

The converse of Theorem~\ref{thm:minimizing-geodesic-unit-speed} holds only
locally.  A regular or piecewise regular curve $\gamma:I\to M$ is
\emph{locally minimizing} if each $t_0\in I$ has a neighborhood $I_0\subseteq
I$ such that $\gamma|_{[a,b]}$ is minimizing whenever $a,b\in I_0$ and
$a<b$.

\begin{lemma}[Subsegments of Minimizing Curves]
\label{lem:minimizing-curves-locally-minimizing}
\dcref{ch6:6.6}

Every minimizing admissible curve segment is locally minimizing.
\end{lemma}

If $\exp_p$ maps $B_\varepsilon(0)\subset T_pM$ diffeomorphically onto its
image, that image is an \emph{open geodesic ball}.  If
$\overline B_\varepsilon(0)$ lies in an open set on which $\exp_p$ is a
diffeomorphism, then $\exp_p(\overline B_\varepsilon(0))$ is a
\emph{closed geodesic ball}, and $\exp_p(\partial B_\varepsilon(0))$ is a
\emph{geodesic sphere}.  Every closed geodesic ball is contained in an open
geodesic ball of larger radius.  In normal coordinates centered at $p$, these
sets are the corresponding coordinate balls and spheres.

On a normal neighborhood $U$ of $p$, define in any normal coordinates
\begin{equation}
r(x) = \sqrt{(x^1)^2 + \cdots + (x^n)^2}, \tag{6.4}
\end{equation}
and, on $U\setminus\{p\}$,
\begin{equation}
\partial_r = \frac{x^i}{r(x)} \frac{\partial}{\partial x^i}. \tag{6.5}
\end{equation}

\begin{lemma}[Radial Objects Are Intrinsic]
\label{lem:radial-distance-well-defined}
\uses{prop:uniqueness-normal-coordinates}
\dcref{ch6:6.8}


The radial distance $r$ and radial vector field $\partial_r$ are independent
of the normal coordinates.  Both are smooth on $U\setminus\{p\}$, and $r^2$
is smooth on $U$.
\end{lemma}

\begin{proof}
\sketch
Normal-coordinate transitions have the form
$\widetilde x^i=A^i_jx^j$ with $A\in O(n)$.  Formulas (6.4) and (6.5) are
invariant under these transitions, and their smoothness statements follow
from the same formulas.
\end{proof}

\begin{theorem}[The Gauss Lemma]
\label{thm:gauss-lemma}
\uses{prop:properties-normal-coordinates}
\dcref{ch6:6.9}


Let $U$ be a geodesic ball centered at $p\in M$.  On
$U\setminus\{p\}$, the radial vector field $\partial_r$ has unit length and is
orthogonal to every geodesic sphere centered at $p$.
\end{theorem}

\begin{proof}
\sketch
Use normal coordinates centered at $p$.  For $q\ne p$, set $b=r(q)$ and
$v^i=q^i/b$.  The radial geodesic $\gamma_v(t)=(tv^1,\ldots,tv^n)$ satisfies
$\gamma_v(b)=q$ and $\gamma_v'(b)=\partial_r|_q$.  Since the metric is
Euclidean at $p$, $|v|=1$, so $\gamma_v$ has unit speed and
$|\partial_r|_q=1$.

Let $w\in T_q\Sigma_b$, where
$\Sigma_b=\exp_p(\partial B_b(0))$, and choose
$\sigma:(-\varepsilon,\varepsilon)\to\Sigma_b$ with
$\sigma(0)=q$ and $\sigma'(0)=w$.  Its coordinate functions satisfy
\begin{equation}
(\sigma^1(s))^2 + \cdots + (\sigma^n(s))^2 = b^2. \tag{6.6}
\end{equation}
Define
\[
\Gamma(s,t)=\left(\frac{t}{b}\sigma^1(s),\ldots,
\frac{t}{b}\sigma^n(s)\right),
\qquad 0\leq t\leq b.
\]
Each main curve is a unit-speed radial geodesic.  With
$S=\partial_s\Gamma$ and $T=\partial_t\Gamma$, one has
$S(0,0)=0$, $S(0,b)=w$, and $T(0,b)=\partial_r|_q$.  Moreover,
\begin{align}
\frac{\partial}{\partial t}\langle S,T\rangle
&=\langle D_tS,T\rangle+\langle S,D_tT\rangle \notag\\
&=\langle D_sT,T\rangle+0 \tag{6.7}\\
&=\frac12\frac{\partial}{\partial s}|T|^2=0. \notag
\end{align}
Therefore $\langle w,\partial_r|_q\rangle
=\langle S,T\rangle(0,b)=\langle S,T\rangle(0,0)=0$.
\end{proof}

\begin{corollary}[Gradient of Radial Distance]
\label{cor:gradient-radial-distance}
\uses{prop:gradient-characterization, prob:gradient-characterization}
\dcref{ch6:6.10}


On a geodesic ball $U$ centered at $p$,
$\operatorname{grad}r=\partial_r$ on $U\setminus\{p\}$.
\end{corollary}

\begin{proof}
\sketch
The Gauss lemma makes $\partial_r$ orthogonal to the level sets of $r$ and
gives $|\partial_r|^2=1$.  Directly from (6.4)--(6.5),
$\partial_r(r)=1$.  The gradient characterization in
\cref{prop:gradient-characterization} now
implies $\operatorname{grad}r=\partial_r$.
\end{proof}

\begin{proposition}[Radial Geodesics Are Unique Minimizers]
\label{prop:radial-geodesic-unique-minimizer}
\lean{LeeLib.Ch06.radialGeodesicUniqueMinimizer}
\leanok
\dcref{ch6:6.11}



Let $p\in M$, and suppose $q$ lies in a geodesic ball centered at $p$.  Up to
reparametrization, the radial geodesic from $p$ to $q$ is the unique
minimizing curve in $M$ joining these points.
\end{proposition}

\begin{proof}
\sketch
Choose a geodesic ball $\exp_p(B_\varepsilon(0))$ containing $q$, and write
the unit-speed radial geodesic as
$\gamma(t)=\exp_p(tv)$ for $0\leq t\leq c$, where $|v|=1$ and
$\gamma(c)=q$.  Thus $L_g(\gamma)=c$.

Let $\sigma:[0,b]\to M$ be any unit-speed admissible curve from $p$ to $q$.
Let $a_0$ be its last visit to $p$, and let $b_0\geq a_0$ be its first
subsequent intersection with the geodesic sphere $\Sigma_c$.  On
$[a_0,b_0]$, the gradient identity and Cauchy--Schwarz give
\begin{align}
r(\sigma(b_0))-r(\sigma(a_0))
&=\int_{a_0}^{b_0}\frac{d}{dt}r(\sigma(t))\,dt \notag\\
&=\int_{a_0}^{b_0}dr(\sigma'(t))\,dt \notag\\
&=\int_{a_0}^{b_0}
\langle\operatorname{grad}r|_{\sigma(t)},\sigma'(t)\rangle\,dt \tag{6.8}\\
&\leq\int_{a_0}^{b_0}
|\operatorname{grad}r|_{\sigma(t)}|\,|\sigma'(t)|\,dt \notag\\
&=L_g(\sigma|_{[a_0,b_0]})\leq L_g(\sigma). \notag
\end{align}
The left side is $c$, so $\gamma$ minimizes length.

If $L_g(\sigma)=c$, equality throughout forces
$a_0=0$, $b_0=b=c$, and equality in Cauchy--Schwarz at every regular point.
Thus $\sigma'$ is a positive multiple of $\operatorname{grad}r$; unit speed
gives $\sigma'=\operatorname{grad}r=\partial_r$.  Hence $\sigma$ and $\gamma$
are the same integral curve through $q$, up to parametrization.
\end{proof}

\begin{corollary}[Radial Distance Equals Riemannian Distance]
\label{cor:radial-distance-riemannian-distance}
\lean{LeeLib.Ch06.radialDistance_eq_riemannianDistance}
\uses{prop:radial-geodesic-unique-minimizer}
\dcref{ch6:6.12}




Let $(M,g)$ be connected and $p\in M$.  On every open or closed geodesic ball
centered at $p$, the radial distance (6.4) equals $d_g(p,\mathord\cdot)$.
\end{corollary}

\begin{proof}
\sketch
It suffices to consider an open geodesic ball because each closed one lies in
a larger open one.  For any $x$ in the ball, the radial geodesic is minimizing
by Proposition~\ref{prop:radial-geodesic-unique-minimizer}.  Its velocity is
$\partial_r$, which has unit length in both the Riemannian metric and normal
coordinates, so its length is $r(x)$.
\end{proof}

\begin{corollary}[Geodesic Balls Are Metric Balls]
\label{cor:geodesic-ball-metric-ball}
\lean{LeeLib.Ch06.geodesicBall_eq_metricBall}
\uses{cor:radial-distance-riemannian-distance}
\leanok
\dcref{ch6:6.13}




In a connected Riemannian manifold, every open or closed geodesic ball is the
metric ball of the same center and radius, and every geodesic sphere is the
corresponding metric sphere.
\end{corollary}

\begin{proof}
\sketch
Let $V=\exp_p(\overline B_c(0))$.  By
Corollary~\ref{cor:radial-distance-riemannian-distance}, points of $V$ have
distance at most $c$ from $p$.  If $q\notin V$, every admissible curve from
$p$ to $q$ crosses $S=\exp_p(\partial B_c(0))$.  Its initial segment to $S$
has length at least $c$ by
Corollary~\ref{cor:radial-distance-riemannian-distance}, and its remaining
segment has positive length; hence $d_g(p,q)>c$.  Thus closed geodesic and
metric balls agree.  Taking unions over radii less than $c$ proves the open
case, and subtracting the open ball from the closed ball proves the sphere
case.
\end{proof}

Henceforth write
\[
B_c(p)=\exp_p(B_c(0)),
\qquad
\overline B_c(p)=\exp_p(\overline B_c(0)),
\qquad
S_c(p)=\exp_p(\partial B_c(0))
\]
for the corresponding geodesic, equivalently metric, balls and spheres.
\section{Uniformly Normal Neighborhoods}

A subset $W$ of a Riemannian manifold $(M,g)$ is \emph{uniformly normal} if
there is a $\delta>0$ such that, for every $x\in W$, the set $W$ is contained
in a geodesic ball of radius $\delta$ centered at $x$. Such a set is also called
\emph{uniformly $\delta$-normal}. Every subset of a uniformly
$\delta$-normal set is uniformly $\delta$-normal.

\begin{lemma}[Uniformly Normal Neighborhood Lemma]
\label{lem:uniformly-normal-neighborhood}
\uses{cor:geodesic-ball-metric-ball}
\dcref{ch6:6.14}

For every $p\in M$ and every neighborhood $U$ of $p$, there is a uniformly
normal neighborhood $W$ of $p$ with $W\subseteq U$.
\end{lemma}

\begin{proof}
\sketch Choose a normal coordinate neighborhood $U_0\subseteq U$ of $p$ and
identify $TU_0$ with $U_0\times\mathbb R^n$. On the domain $\mathcal E$ of the
exponential map, define
\[
E(x,v)=(x,\exp_xv).
\]
Its differential at $(p,0)$ has block form
\[
dE_{(p,0)}=
\begin{pmatrix}
\operatorname{Id}&0\\
*&\operatorname{Id}
\end{pmatrix},
\]
so the inverse function theorem gives neighborhoods on which $E$ is a
diffeomorphism. Shrink the source to $W\times B_\varepsilon(0)$. Then, for
every $x\in W$, the restriction of $\exp_x$ to $B_\varepsilon(0)$ is a
diffeomorphism onto its image, with inverse
$y\mapsto\pi_2(E^{-1}(x,y))$.

After shrinking $W$, a uniform coordinate estimate for $g$ implies that every
$B_\varepsilon(0)\subseteq T_xM$ contains the $g_x$-ball of radius
$\delta=\varepsilon/c$. Thus a geodesic ball of radius $\delta$ is defined at
every $x\in W$. Shrink $W$ once more so that its $d_g$-diameter is less than
$\delta$. Then $W$ lies in the metric $\delta$-ball about every one of its
points, and Corollary~\ref{cor:geodesic-ball-metric-ball} identifies each such
metric ball with the corresponding geodesic ball.
\end{proof}

\begin{theorem}[Every Riemannian geodesic is locally minimizing]
\label{thm:locally-minimizing-geodesics}
\uses{prop:radial-geodesic-unique-minimizer, lem:uniformly-normal-neighborhood}
\dcref{ch6:6.15}

Let $(M,g)$ be a Riemannian manifold and let $\gamma\colon I\to M$ be a
geodesic on an open interval. For every $t_0\in I$, there is an interval
$I_0\subseteq I$ containing $t_0$ such that $\gamma|_{[a,b]}$ is minimizing
whenever $a,b\in I_0$ and $a<b$.
\end{theorem}

\begin{proof}
\sketch Let $W$ be a uniformly normal neighborhood of $\gamma(t_0)$ and let
$I_0$ be the component of $\gamma^{-1}(W)$ containing $t_0$. For
$a<b$ in $I_0$, the point $\gamma(b)$ lies in a geodesic ball centered at
$\gamma(a)$. Proposition~\ref{prop:radial-geodesic-unique-minimizer} makes the
radial segment between these points uniquely minimizing. Since
$\gamma|_{[a,b]}$ is a geodesic in the same ball with the same endpoints, it
equals that radial segment.
\end{proof}

\begin{proof}[Another proof of Theorem~\ref{thm:minimizing-geodesic-unit-speed}]
\sketch Let $\gamma\colon[a,b]\to M$ be a minimizing admissible curve. Near
each $t_0\in[a,b]$, choose a connected interval $I_0$ whose image lies in a
uniformly normal neighborhood. For $a_0<b_0$ in $I_0$, the restriction
$\gamma|_{[a_0,b_0]}$ is minimizing and hence equals the unique minimizing
radial geodesic between its endpoints. Thus $\gamma$ satisfies the geodesic
equation near every $t_0$ and is a geodesic.
\end{proof}

For a Riemannian manifold $(M,g)$ without boundary, define
$\operatorname{inj}(p)$ to be the supremum of all $a>0$ such that
\[
\exp_p\colon B_a(0)\subseteq T_pM\longrightarrow M
\]
is a diffeomorphism onto its image; set $\operatorname{inj}(p)=\infty$ if
these radii are unbounded. The \emph{injectivity radius} of $M$ is
\[
\operatorname{inj}(M)=\inf_{p\in M}\operatorname{inj}(p),
\]
which may be zero, positive, or infinite.

\begin{lemma}
\label{lem:compact-injectivity-radius-positive}
\uses{lem:uniformly-normal-neighborhood}
\dcref{ch6:6.16}

If $(M,g)$ is compact, then $\operatorname{inj}(M)>0$.
\end{lemma}

\begin{proof}
\sketch For every $x\in M$, choose a uniformly $\delta(x)$-normal
neighborhood $W_x$ of $x$. Then
$\operatorname{inj}(x')\geq\delta(x)$ for all $x'\in W_x$. A finite subcover
$W_{x_1},\ldots,W_{x_k}$ gives
\[
\operatorname{inj}(M)\geq\min_i\delta(x_i)>0.
\]
Moreover, it is finite because a compact metric space is bounded, whereas a
geodesic ball of radius $c$ contains points whose distances from its center
approach $c$.
\end{proof}

A subset $U\subseteq M$ is \emph{geodesically convex} if every $p,q\in U$ are
joined by a unique minimizing geodesic segment in $M$ and the image of that
segment is contained in $U$.

\begin{theorem}
\label{thm:small-geodesic-ball-convex}
\dcref{ch6:6.17}
For every $p$ in a Riemannian manifold $(M,g)$, there is an
$\varepsilon_0>0$ such that every geodesic ball centered at $p$ with radius at
most $\varepsilon_0$ is geodesically convex.
\end{theorem}

\section{Completeness}

A connected Riemannian manifold is \emph{metrically complete} if every Cauchy
sequence for $d_g$ converges. It is \emph{geodesically complete} if every
maximal geodesic has domain $\mathbb R$.

\begin{lemma}
\label{lem:exponential-map-implies-completeness}
\lean{LeeLib.Ch06.exponentialMapDefined_implies_completeness}
\leanok
\dcref{ch6:6.18}


Let $(M,g)$ be connected. If some $p\in M$ has $\exp_p$ defined on all of
$T_pM$, then:
\begin{enumerate}
\item[(a)] every $q\in M$ is joined to $p$ by a minimizing geodesic segment;
\item[(b)] $M$ is metrically complete.
\end{enumerate}
\end{lemma}

\begin{proof}
\leanok
\sketch Fix $q\in M$. Say that a minimizing geodesic segment
$\gamma\colon[0,b]\to M$ beginning at $p$ \emph{aims at $q$} if
\begin{equation}
d_g(p,q)=d_g(p,\gamma(b))+d_g(\gamma(b),q). \tag{6.9}
\end{equation}
An aiming segment of length $d_g(p,q)$ ends at $q$.

Choose a closed geodesic ball $\overline B_\varepsilon(p)$ not containing
$q$, and let $x\in S_\varepsilon(p)$ minimize $d_g(x,q)$ on that sphere. Let
$\gamma$ be the maximal unit-speed extension of the radial segment from $p$ to
$x$; it is defined on $\mathbb R$ by the hypothesis on $\exp_p$. Every
admissible curve from $p$ to $q$ first meets $S_\varepsilon(p)$ at some point
$x'$, and splitting there gives
\[
L_g(\sigma)\geq d_g(p,x')+d_g(x',q)
\geq d_g(p,x)+d_g(x,q).
\]
Taking the infimum and using the triangle inequality yields
\begin{equation}
d_g(p,q)=d_g(p,x)+d_g(x,q), \tag{6.10}
\end{equation}
so $\gamma|_{[0,\varepsilon]}$ aims at $q$.

Set $T=d_g(p,q)$ and
\[
\mathcal A=\{b\in[0,T]:\gamma|_{[0,b]}\text{ aims at }q\}.
\]
This nonempty set is closed. Let $A=\sup\mathcal A$, so $A\in\mathcal A$.
If $A<T$, put $y=\gamma(A)$ and choose a closed geodesic
$\delta$-ball about $y$ that excludes $q$. Choose
$z\in S_\delta(y)$ minimizing $d_g(z,q)$, and let $\tau$ be the radial
unit-speed segment from $y$ to $z$. The preceding sphere argument gives
\begin{equation}
d_g(z,q)=d_g(y,q)-d_g(y,z)=(T-A)-\delta. \tag{6.11}
\end{equation}
Consequently $d_g(p,z)\geq A+\delta$. The concatenation of
$\gamma|_{[0,A]}$ and $\tau$ has exactly this length and is therefore
minimizing. A minimizing curve has no corner, so uniqueness for geodesics with
given initial data gives $z=\gamma(A+\delta)$. Equation (6.11) then shows that
$A+\delta\in\mathcal A$, a contradiction. Hence $A=T$, and the aiming segment
$\gamma|_{[0,T]}$ joins $p$ to $q$ and is minimizing.

For (b), let $(q_i)$ be Cauchy. By (a), write
\[
q_i=\exp_p(d_iv_i),
\qquad d_i=d_g(p,q_i),qquad |v_i|=1.
\]
The $d_i$ are bounded, so $(d_iv_i)$ has a convergent subsequence in $T_pM$,
say $d_{i_k}v_{i_k}\to v$. Continuity of $\exp_p$ gives
$q_{i_k}\to\exp_pv$; the Cauchy property forces the full sequence to converge
to the same point.
\end{proof}

\begin{theorem}[Hopf--Rinow]
\label{thm:hopf-rinow}
\lean{LeeLib.Ch06.hopfRinow}
\leanok
\dcref{ch6:6.19}


A connected Riemannian manifold is metrically complete if and only if it is
geodesically complete.
\end{theorem}

\begin{proof}
\uses{lem:exponential-map-implies-completeness}
\leanok
\sketch Geodesic completeness makes every $\exp_p$ defined on all of $T_pM$,
so Lemma~\ref{lem:exponential-map-implies-completeness} gives metric
completeness.

Conversely, suppose $M$ is metrically complete and a unit-speed geodesic
$\gamma\colon[0,b)\to M$ has no extension beyond $b$. For any sequence
$t_i\nearrow b$, the points $q_i=\gamma(t_i)$ satisfy
\[
d_g(q_i,q_j)\leq|t_i-t_j|,
\]
so $q_i\to q\in M$. Choose a uniformly $\delta$-normal neighborhood $W$ of
$q$, and take $j$ with $q_j\in W$ and $t_j>b-\delta$. The geodesic $\sigma$
with initial data $(q_j,\gamma'(t_j))$ is defined on $[0,\delta)$ because the
$\delta$-ball centered at $q_j$ is geodesic. Uniqueness makes
$\sigma(t-t_j)$ agree with $\gamma(t)$ on their overlap, extending $\gamma$
to $[0,t_j+\delta)$ with $t_j+\delta>b$, a contradiction.
\end{proof}

For a connected Riemannian manifold, \emph{complete} means either, hence both,
of the equivalent conditions in Hopf--Rinow. For a disconnected manifold it
means geodesically complete, equivalently that each connected component is a
complete metric space.

\begin{corollary}
\label{cor:exponential-map-defined-implies-complete}
\lean{LeeLib.Ch06.exponentialMapDefined_implies_complete}
\uses{lem:exponential-map-implies-completeness, thm:hopf-rinow}
\leanok
\dcref{ch6:6.20}



If $M$ is connected and $\exp_p$ is defined on all of $T_pM$ for some
$p\in M$, then $M$ is complete.
\end{corollary}

\begin{corollary}
\label{cor:complete-implies-minimizing-geodesics}
\uses{thm:hopf-rinow, lem:exponential-map-implies-completeness}
\dcref{ch6:6.21}

Any two points of a complete, connected Riemannian manifold are joined by a
minimizing geodesic segment.
\end{corollary}

\begin{corollary}
\label{cor:compact-implies-maximal-geodesics-defined}
\uses{thm:hopf-rinow}
\dcref{ch6:6.22}

Every maximal geodesic in a compact Riemannian manifold is defined for all
time.
\end{corollary}

\begin{theorem}
\label{thm:completeness-of-covering-map}
\lean{LeeLib.Ch06.completeLocalIsometry}
\uses{prop:lifting-properties-covering, cor:exponential-map-defined-implies-complete, cor:geodesic-ball-metric-ball}
\leanok
\dcref{ch6:6.23}



Let $(\widetilde M,\widetilde g)$ and $(M,g)$ be connected Riemannian
manifolds. If $\widetilde M$ is complete and
$\pi\colon\widetilde M\to M$ is a local isometry, then $M$ is complete and
$\pi$ is a Riemannian covering map.
\end{theorem}

\begin{proof}
\sketch Fix $\widetilde p\in\widetilde M$, put $p=\pi(\widetilde p)$, and let
$\gamma$ be a geodesic in $M$ with initial point $p$. If
$v=\gamma'(0)$ and
$\widetilde v=(d\pi_{\widetilde p})^{-1}v$, completeness gives a geodesic
$\widetilde\gamma\colon\mathbb R\to\widetilde M$ with initial data
$(\widetilde p,\widetilde v)$. Local isometries preserve geodesics, so
uniqueness gives $\pi\circ\widetilde\gamma=\gamma$ wherever $\gamma$ is
defined. Thus every geodesic from $p$ extends to all of $\mathbb R$, and
Corollary~\ref{cor:exponential-map-defined-implies-complete} makes $M$
complete.

For any $q\in M$, completeness and connectedness give a minimizing unit-speed
geodesic $\gamma$ from $p$ to $q$. Its lift from $\widetilde p$ satisfies
\[
\pi(\widetilde\gamma(r))=q,
\qquad r=d_g(p,q),
\]
so $\pi$ is surjective.

Let $U=B_\varepsilon(p)$ be a geodesic ball and write
$\pi^{-1}(p)=\{\widetilde p_\alpha\}_{\alpha\in A}$. Let
$\widetilde U_\alpha$ be the metric $\varepsilon$-ball about
$\widetilde p_\alpha$. If $\alpha\ne\beta$, a minimizing segment between
$\widetilde p_\alpha$ and $\widetilde p_\beta$ projects to a geodesic loop at
$p$. It must leave and reenter $U$, because geodesics through $p$ contained in
$U$ are radial, and therefore has length at least $2\varepsilon$. Hence
\[
d_{\widetilde g}(\widetilde p_\alpha,\widetilde p_\beta)
\geq2\varepsilon,
\qquad
\widetilde U_\alpha\cap\widetilde U_\beta=\varnothing.
\]

A segment of length less than $\varepsilon$ from $\widetilde p_\alpha$ shows
$\widetilde U_\alpha\subseteq\pi^{-1}(U)$. Conversely, for
$\widetilde q\in\pi^{-1}(U)$, lift the minimizing radial segment from
$q=\pi(\widetilde q)$ to $p$, starting at $\widetilde q$. Its endpoint is
some $\widetilde p_\alpha$, and its length is
$d_g(q,p)<\varepsilon$, so $\widetilde q\in\widetilde U_\alpha$. Thus
\[
\pi^{-1}(U)=\coprod_{\alpha\in A}\widetilde U_\alpha.
\]
For each $\alpha$, the restriction
$\pi\colon\widetilde U_\alpha\to U$ is a local diffeomorphism whose inverse
lifts radial geodesics from $p$ starting at $\widetilde p_\alpha$; it is
therefore a diffeomorphism. Hence $U$ is evenly covered.
\end{proof}

\begin{corollary}
\label{cor:complete-iff-covering-complete}
\lean{LeeLib.Ch06.complete_iff_of_riemannianCovering}
\uses{thm:completeness-of-covering-map}
\leanok
\dcref{ch6:6.24}



Let $\pi\colon\widetilde M\to M$ be a Riemannian covering map between
connected Riemannian manifolds. Then $M$ is complete if and only if
$\widetilde M$ is complete.
\end{corollary}

\begin{proof}
\leanok
\sketch If $\widetilde M$ is complete, apply
Theorem~\ref{thm:completeness-of-covering-map}. Conversely, suppose $M$ is
complete. For arbitrary $\widetilde p\in\widetilde M$ and
$\widetilde v\in T_{\widetilde p}\widetilde M$, the geodesic in $M$ with
initial data
\[
p=\pi(\widetilde p),
\qquad v=d\pi_{\widetilde p}(\widetilde v)
\]
is defined on $\mathbb R$. Its lift starting at $\widetilde p$ is a geodesic
on $\mathbb R$ with initial velocity $\widetilde v$, so $\widetilde M$ is
complete.
\end{proof}

\begin{proposition}
\label{prop:path-homotopy-contains-minimizing-geodesic}
\uses{cor:complete-iff-covering-complete, cor:complete-implies-minimizing-geodesics, prop:lifting-properties-covering}
\dcref{ch6:6.25}

Let $(M,g)$ be complete and connected, and let $p,q\in M$. Every
path-homotopy class of paths from $p$ to $q$ contains a geodesic segment that
has minimum length among all admissible curves in that class.
\end{proposition}

\begin{proof}
\sketch Give the universal cover $\pi\colon\widetilde M\to M$ the pullback
metric. For a representative $\sigma$ of the chosen class, lift it from a
fixed $\widetilde p\in\pi^{-1}(p)$ and let $\widetilde q$ be its endpoint.
The covering space is complete, so there is a minimizing geodesic
$\widetilde\gamma$ from $\widetilde p$ to $\widetilde q$. Its projection
$\gamma=\pi\circ\widetilde\gamma$ belongs to the chosen path-homotopy class.
Every other representative lifts from $\widetilde p$ to the same
$\widetilde q$ by monodromy; its lift is no shorter than
$\widetilde\gamma$. Since $\pi$ is a local isometry, the same length inequality
holds after projection.
\end{proof}

\section{Closed Geodesics}

On a connected Riemannian manifold, a \emph{closed geodesic} is a nonconstant
geodesic segment $\gamma\colon[a,b]\to M$ satisfying
\[
\gamma(a)=\gamma(b),
\qquad
\gamma'(a)=\gamma'(b).
\]

A continuous path $\sigma\colon[0,1]\to M$ is a \emph{loop} if
$\sigma(0)=\sigma(1)$. Loops $\sigma_0$ and $\sigma_1$ are \emph{freely
homotopic} if there is a homotopy $H\colon[0,1]^2\to M$ such that
\begin{align}
H(s,0)=\sigma_0(s),\quad H(s,1)=\sigma_1(s)
&\quad\text{for all }s\in[0,1], \notag\\
H(0,t)=H(1,t)&\quad\text{for all }t\in[0,1]. \tag{6.12}
\end{align}
This is an equivalence relation. Its classes are the \emph{free homotopy
classes}; the class of constant loops is the \emph{trivial free homotopy
class}.

\begin{proposition}[Existence of Closed Geodesics]
\label{prop:existence-closed-geodesics}
\dcref{ch6:6.28}
Let $(M,g)$ be compact and connected. Every nontrivial free homotopy class is
represented by a closed geodesic of minimum length among all admissible loops
in that class.
\end{proposition}

Every compact Riemannian manifold admits a closed geodesic; see
\cite{Jos17} and \cite{Kli95}.
\section{Distance Functions}

For a connected Riemannian manifold $(M,g)$ and $S\subseteq M$, define
\[
d_g(x,S)=\inf\{d_g(x,p):p\in S\}.
\]

\begin{lemma}
\label{lem:distance-triangle-inequality-and-continuity}
\dcref{ch6:6.29}

Let $(M,g)$ be connected and Riemannian, and let $S\subseteq M$.
\begin{enumerate}[label=(\alph*)]
\item For all $x,y\in M$,
$d_g(x,S)\leq d_g(x,y)+d_g(y,S)$.
\item The function $x\mapsto d_g(x,S)$ is continuous on $M$.
\end{enumerate}
\end{lemma}

\begin{theorem}
\label{thm:distance-to-subset-unit-gradient}
\uses{lem:distance-triangle-inequality-and-continuity}
\dcref{ch6:6.31}


Let $(M,g)$ be a connected Riemannian manifold, let $S\subseteq M$, and set
$f(x)=d_g(x,S)$. If $f$ is continuously differentiable on an open set
$U\subseteq M\setminus S$, then
$|\operatorname{grad}f|=1$ on $U$.
\end{theorem}

\begin{proof}
\sketch Fix $x\in U$. For a unit vector $v\in T_xM$, let $\gamma$ be the
unit-speed geodesic with initial data $(x,v)$. For sufficiently small $t>0$,
the distance triangle inequality gives
$f(\gamma(t))\leq t+f(x)$, hence
\[
df_x(v)=\left.\frac d{dt}\right|_{t=0}f(\gamma(t))\leq1.
\]
Taking $v=\operatorname{grad}f|_x/|\operatorname{grad}f|_x$ when the gradient
is nonzero yields $|\operatorname{grad}f|_x\leq1$.

If instead $|\operatorname{grad}f|_x<1$, choose $\delta,\varepsilon>0$ such
that $\overline B_\varepsilon(x)\subseteq U$ and
$|\operatorname{grad}f|\leq1-\delta$ there. Choose
$0<c<\varepsilon\delta$ and an arc-length-parametrized admissible curve
$\alpha:[0,b]\to M$ from $x$ to $S$ with
$b<d_g(x,S)+c$. Since $\overline B_\varepsilon(x)\cap S=\varnothing$,
$b>\varepsilon$, and
\begin{equation}
d_g(\alpha(\varepsilon),S)
\leq b-\varepsilon<d_g(x,S)+c-\varepsilon.
\tag{6.13}
\end{equation}
On $[0,\varepsilon]$,
\[
\left|\frac d{dt}f(\alpha(t))\right|
\leq|\operatorname{grad}f|_{\alpha(t)}|\alpha'(t)|\leq1-\delta,
\]
so
\begin{equation}
d_g(\alpha(\varepsilon),S)
\geq d_g(x,S)-(1-\delta)\varepsilon.
\tag{6.14}
\end{equation}
Together, (6.13) and (6.14) imply $c>\varepsilon\delta$, a contradiction.
\end{proof}

A \emph{local distance function} on an open set $U$ is a continuously
differentiable function $f:U\to\mathbb R$ satisfying
$|\operatorname{grad}f|_g=1$ on $U$.

\begin{theorem}
\label{thm:local-distance-geodesic-integral-curve}
\dcref{ch6:6.32}

If $f$ is a smooth local distance function on an open subset $U$ of a
Riemannian manifold, then
\[
\nabla_{\operatorname{grad}f}(\operatorname{grad}f)=0.
\]
Consequently, every integral curve of $\operatorname{grad}f$ is a unit-speed
geodesic.
\end{theorem}

\begin{proof}
\sketch Put $F=\operatorname{grad}f$. For every vector field $W$,
\begin{equation}
Wf=\langle F,W\rangle,
\tag{6.15}
\end{equation}
and
\begin{equation}
Ff=|F|^2=1.
\tag{6.16}
\end{equation}
Metric compatibility and symmetry give
\begin{align*}
\langle W,\nabla_FF\rangle
&=F\langle W,F\rangle-\langle\nabla_FW,F\rangle\\
&=FWf-[F,W]f-\tfrac12W|F|^2
=W(Ff)-\tfrac12W|F|^2=0.
\end{align*}
Thus $\nabla_FF=0$. An integral curve $\gamma$ satisfies
$D_t\gamma'=\nabla_FF|_\gamma=0$ and $|\gamma'|=|F|=1$.
\end{proof}

\begin{lemma}
\label{lem:local-distance-curve-arc-length}
\dcref{ch6:6.33}

Let $(M,g)$ be Riemannian, let $K\subseteq M$, and let $f:K\to\mathbb R$ be
continuous. Suppose $f|_W$ is a smooth local distance function on an open set
$W\subseteq K$. If an admissible curve
$\sigma:[a_0,b_0]\to K$ satisfies
$\sigma((a_0,b_0))\subseteq W$, then
\[
L_g(\sigma)\geq|f(\sigma(b_0))-f(\sigma(a_0))|.
\]
\end{lemma}

\begin{proof}
\sketch Along the interior of $\sigma$,
$|(f\circ\sigma)'|\leq|\operatorname{grad}f|\,|\sigma'|=|\sigma'|$.
Integrate and use continuity at the endpoints.
\end{proof}

\begin{theorem}
\label{thm:smooth-local-distance-function}
\uses{prop:radial-geodesic-unique-minimizer, thm:local-distance-geodesic-integral-curve, lem:local-distance-curve-arc-length}
\dcref{ch6:6.34}


Let $(M,g)$ be Riemannian, let $U\subseteq M$ be open, let $S\subseteq U$,
and let $f:U\to[0,\infty)$ be continuous with $f^{-1}(0)=S$. If $f$ is a
smooth local distance function on $U\setminus S$, then some neighborhood
$U_0\subseteq U$ of $S$ satisfies
\[
f(x)=d_g(x,S)
\qquad(x\in U_0).
\]
\end{theorem}

\begin{proof}
\sketch For each $p\in S$, choose $\varepsilon_p,\delta_p>0$ such that
$B_{\varepsilon_p}(p)$ is uniformly $\delta_p$-normal and
$B_{\delta_p}(p)\subseteq U$. Set
$U_0=\bigcup_{p\in S}B_{\varepsilon_p}(p)$.

Fix $x\in U_0\setminus S$, choose $p\in S$ with
$x\in B_{\varepsilon_p}(p)$, and write $c=f(x)$. Applying
Lemma~\ref{lem:local-distance-curve-arc-length} to the radial segment from
$p$ to $x$ shows $c<\varepsilon_p$. The geodesic from $x$ with initial
velocity $-\operatorname{grad}f|_x$ follows the negative gradient flow while
it remains outside $S$, and there
\[
\frac d{dt}f(\gamma(t))=-|\operatorname{grad}f|^2=-1.
\]
Thus $f(\gamma(t))=c-t$ for $t<c$, continuity gives
$\gamma(c)\in S$, and $d_g(x,S)\leq c$.

Conversely, let $\alpha$ be any admissible curve from $x$ to $S$. If its
image lies in $U$, stop it at its first intersection with $S$ and apply the
lemma to obtain $L_g(\alpha)\geq c$. If it leaves $U$, the uniform normal-ball
choice forces it first to travel distance at least $\varepsilon_p>c$. Hence
every such curve has length at least $c$, so $d_g(x,S)=c$.
\end{proof}

\begin{corollary}
\label{cor:local-distance-smooth-level-sets}
\uses{thm:smooth-local-distance-function}
\dcref{ch6:6.35}


Let $f$ be a smooth local distance function on an open subset $U$ of a
Riemannian manifold. If $c\in\mathbb R$ and
$S=f^{-1}(c)\neq\varnothing$, then some neighborhood $U_0$ of $S$ in $U$
satisfies
\[
|f(x)-c|=d_g(x,S)
\qquad(x\in U_0).
\]
\end{corollary}

\section{Distance Functions for Embedded Submanifolds}

Let $P$ be an embedded $k$-dimensional submanifold of a Riemannian
$n$-manifold $(M,g)$ without boundary, and let $NP$ be its normal bundle.

\begin{proposition}
\label{prop:fermi-coordinates-radial-function}
\dcref{ch6:6.37}

Let $U$ be a normal neighborhood of $P$. There are unique objects consisting
of a continuous function $r:U\to[0,\infty)$ and a smooth vector field
$\partial_r$ on $U\setminus P$ such that in every Fermi coordinate system
$(x^1,\ldots,x^k,v^1,\ldots,v^{n-k})$,
\begin{equation}
r(x,v)=\sqrt{(v^1)^2+\cdots+(v^{n-k})^2},
\tag{6.17}
\end{equation}
and
\begin{equation}
\partial_r
=\frac{v^1}{r(x,v)}\frac{\partial}{\partial v^1}
+\cdots+
\frac{v^{n-k}}{r(x,v)}\frac{\partial}{\partial v^{n-k}}.
\tag{6.18}
\end{equation}
Moreover, $r$ is smooth on $U\setminus P$, and $r^2$ is smooth on $U$.
\end{proposition}

\begin{proof}
\sketch Let $V\subseteq NP$ be mapped diffeomorphically onto $U$ by the
normal exponential map $E$, put $\rho(p,v)=|v|_g$, and define
$r=\rho\circ E^{-1}$. An orthonormal local frame of $NP$ gives (6.17), which
also proves the asserted regularity.

For $q=\exp_p(v)\in U\setminus P$, let $\gamma(t)=\exp_p(tv)$ and define
\begin{equation}
\partial_r|_q=\frac1{r(q)}\gamma'(1).
\tag{6.19}
\end{equation}
In Fermi coordinates, $\gamma(t)=(x^1,\ldots,x^k,tv^1,\ldots,tv^{n-k})$,
so (6.19) gives (6.18), independently of the chosen coordinates.
\end{proof}

The objects $r$ and $\partial_r$ are respectively the \emph{radial distance
function} and \emph{radial vector field} for $P$.

\begin{theorem}[Gauss Lemma for Submanifolds]
\label{thm:gauss-lemma-submanifolds}
\uses{prop:fermi-coordinates-radial-function}
\dcref{ch6:6.38}


Under the preceding hypotheses, $\partial_r$ is a unit vector field on
$U\setminus P$ and is orthogonal to every level set of $r$.
\end{theorem}

\begin{proof}
\sketch If $q=\exp_p(v)$ and $\gamma(t)=\exp_p(tv)$, then
$|\gamma'(0)|=|v|=r(q)$ and geodesic speed is constant. Formula (6.19) therefore
implies $|\partial_r|_q=1$.

Let $q=\exp_{p_0}(v_0)$, set $b=|v_0|=r(q)>0$, and take
$w\in T_qr^{-1}(b)$. Choose
$\sigma(s)=\exp_{x(s)}(v(s))$ in $r^{-1}(b)$ with
$\sigma(0)=q$, $\sigma'(0)=w$, and define the geodesic variation
\[
\Gamma(s,t)=\exp_{x(s)}\left(\frac tbv(s)\right),
\qquad 0\leq t\leq b.
\]
For $S=\partial_s\Gamma$ and $T=\partial_t\Gamma$, the Gauss-lemma
calculation gives $\partial_t\langle S,T\rangle=0$. Its endpoint values yield
\[
\langle w,\partial_r|_q\rangle
=\frac1b\langle x'(0),v_0\rangle=0,
\]
because $x'(0)$ is tangent to $P$ and $v_0$ is normal to $P$.
\end{proof}

\begin{corollary}[Assume the hypotheses of Theorem~6.38]
\label{cor:distance-function-local-properties}
\uses{thm:gauss-lemma-submanifolds, thm:small-geodesic-ball-convex, lem:exponential-map-implies-completeness, prop:gradient-characterization, prop:fermi-coordinates-properties, thm:smooth-local-distance-function, prob:gradient-characterization}
\dcref{ch6:6.39}


\begin{enumerate}
\item[(a)] On $U\setminus P$, $\partial_r=\operatorname{grad}r$.
\item[(b)] The function $r$ is a local distance function on $U\setminus P$.
\item[(c)] Every unit-speed geodesic $\gamma:[a,b]\to U$ whose initial
velocity is normal to $P$ agrees with an integral curve of $\partial_r$ on
$(a,b)$.
\item[(d)] Some tubular neighborhood of $P$ satisfies
$d_g(x,P)=r(x)$.
\end{enumerate}
\end{corollary}

\begin{proof}
\sketch Formulas (6.17)--(6.18) give
$\partial_r(r)=1=|\partial_r|^2$, and the Gauss lemma gives that
$\partial_r$ is normal to the level sets of $r$. Thus
$\partial_r=\operatorname{grad}r$ and $r$ is a local distance function.
In Fermi coordinates, a unit-speed geodesic normal to $P$ is
\[
t\longmapsto(x^1,\ldots,x^k,tv^1,\ldots,tv^{n-k}),
\qquad \sum_a(v^a)^2=1,
\]
which is an integral curve of $\partial_r$ wherever $r\neq0$. Finally apply
Theorem~\ref{thm:smooth-local-distance-function} and restrict its neighborhood
to a tubular neighborhood of $P$.
\end{proof}

\begin{theorem}
\label{thm:tubular-neighborhood-metric-distance}
\uses{prop:fermi-coordinates-radial-function, lem:local-distance-curve-arc-length}
\dcref{ch6:6.40}


Let $(M,g)$ be connected and Riemannian, let $P\subseteq M$ be a compact
submanifold, and let $U_\varepsilon$ be an $\varepsilon$-tubular neighborhood
of $P$. Then
\[
U_\varepsilon=\{q\in M:d_g(q,P)<\varepsilon\},
\]
and $d_g(q,P)=r(q)$ for every $q\in U_\varepsilon$.
\end{theorem}

\begin{proof}
\sketch First extend $r$ continuously to $\overline U_\varepsilon$ by setting
$r=\varepsilon$ on $\partial U_\varepsilon$. Indeed, if
$q_i=\exp_{p_i}(v_i)\to q\in\partial U_\varepsilon$ and a subsequence had
$r(q_i)=|v_i|\to c<\varepsilon$, compactness of $P$ would yield a further
subsequence $(p_i,v_i)\to(p,v)$ with $|v|=c$. Then
$q=\exp_p(v)\in U_\varepsilon$, a contradiction.

Let $W_\varepsilon=\{q:d_g(q,P)<\varepsilon\}$. Any admissible curve from
$P$ to a point outside $U_\varepsilon$ first meets
$\partial U_\varepsilon$; Lemma~\ref{lem:local-distance-curve-arc-length}
shows that its length is at least $\varepsilon$. Hence
$W_\varepsilon\subseteq U_\varepsilon$.

For $q\in U_\varepsilon$, the normal radial geodesic has length $r(q)$, so
$d_g(q,P)\leq r(q)$. Conversely, an admissible curve from $P$ to $q$ that
remains in $U_\varepsilon$ has length at least $r(q)$ by the same lemma,
applied after its last intersection with $P$. A curve that leaves
$U_\varepsilon$ has length at least $\varepsilon>r(q)$. Thus
$d_g(q,P)=r(q)<\varepsilon$, proving the reverse inclusion and both claims.
\end{proof}

\section{Semigeodesic Coordinates}

Smooth coordinates $(x^1,\ldots,x^n)$ on an open subset of a Riemannian
$n$-manifold are \emph{semigeodesic} if every $x^n$-coordinate curve
\[
t\longmapsto(x^1,\ldots,x^{n-1},t)
\]
is a unit-speed geodesic orthogonal to each level set of $x^n$. In this
section, Latin indices range from $1$ to $n$, and Greek indices range from
$1$ to $n-1$.

\begin{proposition}[Characterizations of Semigeodesic Coordinates]
\label{prop:characterizations-semigeodetic-coordinates}
\uses{thm:local-distance-geodesic-integral-curve}
\dcref{ch6:6.41}


For smooth coordinates $(x^1,\ldots,x^n)$ on an open subset of a Riemannian
$n$-manifold, the following are equivalent:
\begin{enumerate}
\item[(a)] The coordinates are semigeodesic.
\item[(b)] $|\partial_n|_g=1$ and
$\langle\partial_\alpha,\partial_n\rangle_g=0$ for
$1\leq\alpha\leq n-1$.
\item[(c)] $|dx^n|_g=1$ and
$\langle dx^\alpha,dx^n\rangle_g=0$ for
$1\leq\alpha\leq n-1$.
\item[(d)] $|\operatorname{grad}x^n|_g=1$ and
$\langle\operatorname{grad}x^\alpha,\operatorname{grad}x^n\rangle_g=0$ for
$1\leq\alpha\leq n-1$.
\item[(e)] The function $x^n$ is a local distance function, and
$x^1,\ldots,x^{n-1}$ are constant along the integral curves of
$\operatorname{grad}x^n$.
\item[(f)] $\operatorname{grad}x^n=\partial_n$.
\end{enumerate}
\end{proposition}

\begin{proof}
\sketch Condition (b) says that the matrix of $g$ has block form
$\left(\begin{smallmatrix} *&0\\0&1\end{smallmatrix}\right)$; condition (c)
says the same about its inverse, so they are equivalent. Raising covector
indices gives (c)$\Leftrightarrow$(d). Moreover,
\[
(\operatorname{grad}x^n)(x^\alpha)
=\langle\operatorname{grad}x^\alpha,\operatorname{grad}x^n\rangle,
\]
which proves (d)$\Leftrightarrow$(e), and
$\operatorname{grad}x^n=g^{nj}\partial_j$ proves
(c)$\Leftrightarrow$(f).

Semigeodesicity directly implies (b). Conversely, (b) and (f) make the
$x^n$-coordinate curves integral curves of the unit vector field
$\operatorname{grad}x^n$; they are unit-speed geodesics by
Theorem~\ref{thm:local-distance-geodesic-integral-curve} and are orthogonal to
the $x^n$-level sets by (b).
\end{proof}

\begin{corollary}
\label{cor:semigeodesic-metric-christoffel}
\uses{prop:characterizations-semigeodetic-coordinates, prop:tangential-connection-compatible-metric}
\dcref{ch6:6.42}


Let $(x^i)$ be semigeodesic coordinates on an open subset of a Riemannian
$n$-manifold.
\begin{enumerate}
\item[(a)] The metric is
\[
g=(dx^n)^2+g_{\alpha\beta}(x^1,\ldots,x^n)\,dx^\alpha\,dx^\beta.
\]
\item[(b)] Its Christoffel symbols satisfy
\begin{align*}
\Gamma^n_{nn}&=\Gamma^\alpha_{nn}=\Gamma^n_{\alpha n}
=\Gamma^n_{n\alpha}=0,\\
\Gamma^n_{\alpha\beta}&=-\frac12\partial_ng_{\alpha\beta},\\
\Gamma^\beta_{n\alpha}=\Gamma^\beta_{\alpha n}
&=\frac12g^{\beta\gamma}\partial_ng_{\alpha\gamma},\\
\Gamma^\gamma_{\alpha\beta}&=\widetilde\Gamma^\gamma_{\alpha\beta},
\end{align*}
where, for fixed $x^n$, $\widetilde\Gamma^\gamma_{\alpha\beta}$ are the
Christoffel symbols in $(x^\alpha)$ coordinates of the induced metric on the
level set $x^n=\text{constant}$.
\end{enumerate}
\end{corollary}

\begin{proof}
\sketch Proposition~\ref{prop:characterizations-semigeodetic-coordinates}(b)
gives $g_{nn}=1$ and $g_{\alpha n}=0$, which proves (a). Substitute these
identities into formula (5.8) for the Christoffel symbols to obtain (b).
\end{proof}

Semigeodesic coordinates can be constructed from a smooth local distance
function $r$: set $x^n=r$, choose coordinates $x^1,\ldots,x^{n-1}$ on a level
set, and extend them constantly along the integral curves of
$\operatorname{grad}r$.

\begin{example}[Fermi Coordinates for a Hypersurface]
\label{ex:fermi-coordinates-hypersurface}
\uses{thm:tubular-neighborhood, thm:small-geodesic-ball-convex, cor:distance-function-local-properties, lem:exponential-map-implies-completeness, prop:characterizations-semigeodetic-coordinates}
\dcref{ch6:6.43}


Let $P$ be an embedded hypersurface of $(M,g)$, and let
$(x^1,\ldots,x^{n-1},v)$ be Fermi coordinates on $U$. Then
\[
r(x,v)=|v|,
\qquad
\partial_r=\pm\partial_v=\operatorname{grad}|v|
\quad\text{on }U\setminus P.
\]
The coordinate $v$ is a local distance function on $U\setminus P$, and
$|v|=d_g(\,\cdot\,,P)$ on a neighborhood of $P$. On both components of
$U\setminus P$, the displayed identities give
$\operatorname{grad}v=\partial_v$; smoothness extends this equality across
$P$. Thus Fermi coordinates for a hypersurface are semigeodesic.
\end{example}

\begin{example}[Boundary Normal Coordinates]
\label{ex:boundary-normal-coordinates}
\uses{prop:double-manifold}
\dcref{ch6:6.44}


Let $(M,g)$ be a smooth Riemannian manifold with nonempty boundary. Embed $M$
in its double $\widetilde M$, extend $g$ smoothly, and choose Fermi coordinates
$(x^1,\ldots,x^{n-1},v)$ for $\partial M$ in $\widetilde M$. After replacing
$v$ by $-v$ if necessary, require $v>0$ on $\operatorname{Int}M$. Their
restrictions to $M$ are semigeodesic boundary coordinates, called
\emph{boundary normal coordinates}.
\end{example}

\begin{example}[Polar Normal Coordinates]
\label{ex:polar-normal-coordinates}
\dcref{ch6:6.45}

Let $\psi:\widetilde U\to U\subseteq S^{n-1}$ be a smooth local
parametrization and define
\[
\Psi:\widetilde U\times\mathbb R^+\to\mathbb R^n,
\qquad
\Psi(\theta^1,\ldots,\theta^{n-1},r)=r\psi(\theta^1,\ldots,\theta^{n-1}).
\]
Its differential is everywhere injective, so $\Theta=\Psi^{-1}$ gives polar
coordinates on
$\Psi(\widetilde U\times\mathbb R^+)\subseteq\mathbb R^n\setminus\{0\}$,
with last coordinate $r(x)=|x|$.

If $\varphi$ is a normal coordinate chart on a normal neighborhood $V$ of
$p$ in a Riemannian manifold, then $\Theta\circ\varphi$ defines
\emph{polar normal coordinates} on an open subset of $V\setminus\{p\}$. The
last coordinate is radial distance, while the angular coordinates are constant
along its gradient curves; hence these coordinates are semigeodesic.
\end{example}

\begin{example}[Polar Fermi Coordinates]
\label{ex:polar-fermi-coordinates}
\uses{prop:fermi-coordinates-properties}
\dcref{ch6:6.46}


Let $P\subseteq(M,g)$ be embedded, and let
$\varphi=(x^1,\ldots,x^k,v^1,\ldots,v^{n-k})$ be Fermi coordinates on a
neighborhood $U_0$ of $p\in P$. Applying any polar coordinate map $\Theta$
for $\mathbb R^{n-k}$ to the normal variables gives
\[
(\operatorname{Id}_{\mathbb R^k}\times\Theta)\circ\varphi
\]
on an open subset of $U_0\setminus P$, with values in
$\mathbb R^k\times\mathbb R^{n-k-1}\times\mathbb R^+$. Written as
$(x^1,\ldots,x^k,\theta^1,\ldots,\theta^{n-k-1},r)$, every $r$-coordinate
curve is a unit-speed normal geodesic. These semigeodesic coordinates are
\emph{polar Fermi coordinates}; polar normal coordinates are the case
$P=\{p\}$.
\end{example}
\section{Further Results}

\begin{proposition}
\label{prop:round-metric-distance-sphere}
\dcref{ch6:6-2}

Let $n\in\mathbb Z_{>0}$ and $R>0$.
\begin{enumerate}
\item[(a)] For $p,q\in S^n(R)$ with the round metric,
\[
d_{\bar g_R}(p,q)=R\arccos\frac{\langle p,q\rangle}{R^2},
\]
where the inner product is Euclidean.
\item[(b)] The metric space $(S^n(R),d_{\bar g_R})$ has diameter $\pi R$.
\end{enumerate}
\end{proposition}

\begin{exercise}
\label{prob:round-metric-distance-sphere}
\dcref{ch6:6-2}

Let $n$ be a positive integer and $R$ a positive real number.
\begin{enumerate}
\item[(a)] Prove that the Riemannian distance between any two points $p, q$ in $S^n(R)$ with the round metric is given by
\[
d_{\bar{g}_R}(p,q) = R \arccos \frac{\langle p,q\rangle}{R^2},
\]
where $\langle\cdot,\cdot\rangle$ is the Euclidean inner product on $\mathbb{R}^{n+1}$.
\item[(b)] Prove that the metric space $(S^n(R), d_{\bar{g}_R})$ has diameter $\pi R$.
\end{enumerate}
(Used on pp. 39, 359.)
\end{exercise}


\begin{proposition}
\label{prop:isometry-group-acts-properly}
\dcref{ch6:6-28}

Let $(M,g)$ be complete, connected, and $n$-dimensional, and assume its
isometry group is a Lie group acting smoothly.  If $G$ is a closed subgroup of
$\operatorname{Iso}(M)$, then its action on $M$ is proper.
\end{proposition}

\begin{exercise}
\label{prob:isometry-group-acts-properly}
\dcref{ch6:6-28}

Suppose $(M, g)$ is a complete, connected Riemannian $n$-manifold whose isometry group is a Lie group acting smoothly on $M$. (The Myers--Steenrod theorem shows that this is always the case when $M$ is connected, but we have not proved this.) Prove that if $G$ is a closed subgroup of $\operatorname{Iso}(M)$, then $G$ acts properly on $M$, using the following outline. Suppose $(\varphi_m)$ is a sequence in $G$ and $(p_m)$ is a sequence in $M$ such that $p_m \to p_0$ and $\varphi_m(p_m) \to q_0$.
\begin{enumerate}
\item[(a)] Prove that $\varphi_m(p_0) \to q_0$.
\item[(b)] Let $B_r(p_0)$ be a geodesic ball centered at $p_0$; let $0 < \varepsilon < r$; let $\{b_1, \ldots, b_n\}$ be an orthonormal basis for $T_{p_0}M$; and let $p_i = \exp_{p_0}(\varepsilon b_i)$ for $i = 1, \ldots, n$. Prove that there exist a linear isometry $A : T_{p_0}M \to T_{q_0}M$ and a subsequence $(\varphi_{m_j})$ such that $\varphi_{m_j}(p_i) \to \exp_{q_0}(\varepsilon A b_i)$ for $i = 1, \ldots, n$. [Hint: Use Prop. 5.20.]
\item[(c)] Prove that there is an isometry $\varphi \in G$ such that $\varphi_{m_j} \to \varphi$ pointwise, meaning that $\varphi_{m_j}(x) \to \varphi(x)$ for every $x \in M$.
\item[(d)] Prove that $\varphi_{m_j} \to \varphi$ in the topology of $G$. [Hint: Define a map $F: G \to M^{n+1}$ by $F(\psi) = (\psi(p_0), \psi(p_1), \ldots, \psi(p_n))$, where $p_1, \ldots, p_n$ are the points introduced in part (b). Show that $F$ is continuous, injective, and closed, and therefore is a homeomorphism onto its image.]
\end{enumerate}
(Used on p. 349.)
\end{exercise}


\begin{exercise}
\label{exer:closed-geodesic-periodic-extension}
\dcref{ch6:6.26}

Show that a geodesic segment is closed if and only if it extends to a periodic geodesic defined on all of $\mathbb{R}$.
\end{exercise}

\begin{exercise}
\label{exer:prove-distance-lemma}
\dcref{ch6:6.30}
\uses{lem:distance-triangle-inequality-and-continuity}

Prove the preceding lemma.
\end{exercise}

\begin{exercise}
\label{exer:prove-local-distance-corollary}
\dcref{ch6:6.36}
\uses{cor:local-distance-smooth-level-sets}

Prove the preceding corollary.
\end{exercise}

\begin{exercise}
\label{exer:prove-minimizing-lemma}
\dcref{ch6:6.7}
\uses{lem:minimizing-curves-locally-minimizing}

Prove the preceding lemma.
\end{exercise}

\begin{exercise}
\label{exer:trivial-free-homotopy-identity}
\dcref{ch6:6.27}

Given a connected manifold $M$ and a point $x \in M$, show that a loop based at $x$ represents the trivial free homotopy class if and only if it represents the identity element of $\pi_1(M, x)$.
\end{exercise}

\begin{exercise}
\label{prob:closed-geodesic-free-homotopy}
\dcref{ch6:6-17}
\uses{prop:existence-closed-geodesics,prop:path-homotopy-contains-minimizing-geodesic}

Prove Proposition~\ref{prop:existence-closed-geodesics} (existence of a closed geodesic in a free homotopy class). [Hint: Use Prop.~\ref{prop:path-homotopy-contains-minimizing-geodesic} to show that the given free homotopy class is represented by a \emph{geodesic loop}, i.e., a geodesic whose starting and ending points are the same. Show that the lengths of such loops have a positive greatest lower bound; then choose a sequence of geodesic loops whose lengths approach that lower bound, and show that a subsequence converges uniformly to a closed geodesic.]
\end{exercise}

\begin{exercise}
\label{prob:compact-complete-characterization}
\dcref{ch6:6-14}

Let $(M, g)$ be a connected Riemannian manifold.
\begin{enumerate}
\item[(a)] Show that $M$ is complete if and only if the compact subsets of $M$ are exactly the closed and bounded ones.
\item[(b)] Show that $M$ is compact if and only if it is complete and bounded.
\end{enumerate}
\end{exercise}

\begin{exercise}
\label{prob:completeness-divergence}
\dcref{ch6:6-10}

A curve $\gamma : [0, b) \to M$ (with $0 < b \leq \infty$) is said to \emph{diverge to infinity} if for every compact set $K \subseteq M$, there is a time $T \in [0, b)$ such that $\gamma(t) \notin K$ for $t > T$. (For those who are familiar with one-point compactifications, this means that $\gamma(t)$ converges to the ``point at infinity'' in the one-point compactification of $M$ as $t \nearrow b$.) Prove that a connected Riemannian manifold is complete if and only if every regular curve that diverges to infinity has infinite length. (The length of a curve whose domain is not compact is just the supremum of the lengths of its restrictions to compact subintervals.)
\end{exercise}

\begin{exercise}
\label{prob:connected-riemannian-1-manifold-structure}
\dcref{ch6:6-1}

Suppose $M$ is a nonempty connected Riemannian 1-manifold. Show that if $M$ is noncompact, then it is isometric to an open interval in $\mathbb{R}$ with the Euclidean metric, while if it is compact, it is isometric to a circle $S^1(R) = \{x \in \mathbb{R}^2 : |x| = R\}$ with its induced metric for some $R > 0$, using the following steps.
\begin{enumerate}
\item[(a)] Let $\gamma : I \to M$ be any maximal unit-speed geodesic. Show that its image is open and closed, and therefore $\gamma$ is surjective.
\item[(b)] Show that if $\gamma$ is injective, then it is an isometry between $I$ with its Euclidean metric and $M$.
\item[(c)] Now suppose $\gamma(t_1) = \gamma(t_2)$ for some $t_1 \ne t_2$. In case $\gamma'(t_1) = \gamma'(t_2)$, show that $\gamma$ is periodic, and descends to a global isometry from an appropriate circle to $M$.
\item[(d)] It remains only to rule out the case $\gamma(t_1) = \gamma(t_2)$ and $\gamma'(t_1) = -\gamma'(t_2)$. If this occurs, let $t_0 = (t_1 + t_2)/2$, and define geodesics $\alpha$ and $\beta$ by
\[
\alpha(t) = \gamma(t_0 + t), \quad \beta(t) = \gamma(t_0 - t).
\]
Use uniqueness of geodesics to conclude that $\alpha \equiv \beta$ on their common domain, and show that this contradicts the fact that $\gamma$ is injective on some neighborhood of $t_0$.
\end{enumerate}
\end{exercise}

\begin{exercise}[Convexity radius is continuous]
\label{prob:convexity-radius-continuous}
\dcref{ch6:6-6}

Suppose $(M, g)$ is a Riemannian manifold. For each $x \in M$, define the \emph{convexity radius of $M$ at $x$}, denoted $\operatorname{conv}(x)$, to be the supremum of all $\varepsilon > 0$ such that there is a geodesically convex geodesic ball of radius $\varepsilon$ centered at $x$. Show that $\operatorname{conv}(x)$ is a continuous function of $x$.
\end{exercise}

\begin{exercise}
\label{prob:distance-complete-metric}
\dcref{ch6:6-9}
\uses{prob:distance-disconnected-riemannian}

Suppose $(M, g)$ is a (not necessarily connected) Riemannian manifold. In Problem 2-30, you were asked to show that there is a distance function $d : M \times M \to \mathbb{R}$ that induces the given topology and restricts to the Riemannian distance on each connected component of $M$. Show that if each component of $M$ is geodesically complete, then $d$ can be chosen so that $(M, d)$ is a complete metric space. Show also that if $M$ has infinitely many components, then $d$ can be chosen so that $(M, d)$ is not complete.
\end{exercise}

\begin{exercise}[Metric and Riemannian distance]
\label{prob:distance-determines-metric}
\dcref{ch6:6-4}

In Chapter 2, we started with a Riemannian metric and used it to define the Riemannian distance function. This problem shows how to go back the other way: the distance function determines the Riemannian metric. Let $(M, g)$ be a connected Riemannian manifold.
\begin{enumerate}
\item[(a)] Show that if $\gamma: (-\varepsilon, \varepsilon) \to M$ is any smooth curve, then
\[
|\gamma'(0)|_g = \lim_{t \searrow 0} \frac{d_g(\gamma(0), \gamma(t))}{t}.
\]
\item[(b)] Show that if $g$ and $\tilde{g}$ are two Riemannian metrics on $M$ such that $d_g(p,q) = d_{\tilde{g}}(p,q)$ for all $p,q \in M$, then $g = \tilde{g}$.
\end{enumerate}
\end{exercise}

\begin{exercise}
\label{prob:distance-function-submersion}
\dcref{ch6:6-25}

Suppose $(M, g)$ is a Riemannian manifold and $f : M \to \mathbb{R}$ is a smooth local distance function. Prove that $f$ is a Riemannian submersion.
\end{exercise}

\begin{exercise}
\label{prob:energy-functional}
\dcref{ch6:6-23}

In some treatments of Riemannian geometry, instead of minimizing the length functional, one considers the following \emph{energy functional} for an admissible curve $\gamma: [a,b] \to M$:
\[
E(\gamma) = \frac{1}{2} \int_a^b |\gamma'(t)|^2 \, dt.
\]
(Note that $E(\gamma)$ is \emph{not} independent of parametrization.)
\begin{enumerate}
\item[(a)] Prove that an admissible curve is a critical point for $E$ (with respect to proper variations) if and only if it is a geodesic (which means, in particular, that it has constant speed).
\item[(b)] Prove that if $\gamma$ is an admissible curve that minimizes energy among admissible curves with the same endpoints, then it also minimizes length.
\item[(c)] Prove that if $\gamma$ is an admissible curve that minimizes length among admissible curves with the same endpoints, then it minimizes energy if and only if it has constant speed.
\end{enumerate}
[Remark: For our limited purposes, it is easier and more straightforward to use the length functional. But because the energy functional does not involve the square root function, and its critical points automatically have constant-speed parametrizations, it is sometimes more useful for proving the existence of geodesics with certain properties.]
\end{exercise}

\begin{exercise}
\label{prob:extendible}
\dcref{ch6:6-20}

A connected Riemannian manifold is said to be \emph{extendible} if it is isometric to a proper open subset of a larger connected Riemannian manifold.
\begin{enumerate}
\item[(a)] Show that every complete, connected Riemannian manifold is nonextendible.
\item[(b)] Show that the converse is false by giving an example of a nonextendible connected Riemannian manifold that is not complete.
\end{enumerate}
\end{exercise}

\begin{exercise}[Geodesic balls are geodesically convex]
\label{prob:geodesic-ball-convexity}
\dcref{ch6:6-5}
\uses{thm:small-geodesic-ball-convex}

Prove Theorem~\ref{thm:small-geodesic-ball-convex} (sufficiently small geodesic balls are geodesically convex) as follows.
\begin{enumerate}
\item[(a)] Let $(M,g)$ be a Riemannian manifold, let $p \in M$ be fixed, and let $W$ be a uniformly normal neighborhood of $p$. For $\varepsilon > 0$ small enough that $B_{2\varepsilon}(p) \subseteq W$, define a subset $W_\varepsilon \subseteq TM \times \mathbb{R}$ by
\[
W_\varepsilon = \{(q, v, t) \in TM \times \mathbb{R} : q \in B_\varepsilon(p), v \in T_q M, |v| = 1, |t| < 2\varepsilon\}.
\]
Define $f: W_\varepsilon \to \mathbb{R}$ by
\[
f(q, v, t) = d_g(p, \exp_q(tv))^2.
\]
Show that $f$ is smooth. [Hint: Use normal coordinates centered at $p$.]
\item[(b)] Show that if $\varepsilon$ is chosen small enough, then $\frac{\partial^2 f}{\partial t^2} > 0$ on $W_\varepsilon$. [Hint: Compute $f(p, v, t)$ explicitly and use continuity. Be careful to verify that $\varepsilon$ can be chosen independently of $v$.]
\item[(c)] Suppose $\varepsilon$ is chosen as in (b). Show that if $q_1, q_2 \in B_\varepsilon(p)$ and $\gamma$ is a minimizing geodesic segment from $q_1$ to $q_2$, then $d_g(p, \gamma(t))$ attains its maximum at one of the endpoints of $\gamma$.
\item[(d)] Show that $B_\varepsilon(p)$ is geodesically convex.
\end{enumerate}
\end{exercise}

\begin{exercise}
\label{prob:geodesic-completeness}
\dcref{ch6:6-12}

Let $(M, g)$ be a connected Riemannian manifold.
\begin{enumerate}
\item[(a)] Suppose there exists $\delta > 0$ such that for each $p \in M$, every maximal unit-speed geodesic starting at $p$ is defined at least on an interval of the form $(-\delta, \delta)$. Prove that $M$ is complete.
\item[(b)] Prove that if $M$ has positive or infinite injectivity radius, then it is complete.
\item[(c)] Prove that if $M$ is homogeneous, then it is complete.
\item[(d)] Give an example of a complete, connected Riemannian manifold that has zero injectivity radius.
\end{enumerate}
\end{exercise}

\begin{exercise}
\label{prob:geodesic-to-closed-subset}
\dcref{ch6:6-26}

Suppose $(M, g)$ is a complete, connected Riemannian manifold and $S \subseteq M$ is a closed subset.
\begin{enumerate}
\item[(a)] Show that for every $p \in M \setminus S$, there is a geodesic segment from $p$ to a point of $S$ whose length is equal to $d_g(p, S)$.
\item[(b)] In case $S$ is a properly embedded submanifold of $M$, show that every geodesic segment satisfying the conclusion of (a) intersects $S$ orthogonally.
\item[(c)] Give a counterexample to (a) if $S$ is not closed.
\end{enumerate}
\end{exercise}

\begin{exercise}
\label{prob:inj-radius-diffeomorphism}
\dcref{ch6:6-16}

Suppose $(M, g)$ is a complete, connected Riemannian manifold with positive or infinite injectivity radius.
\begin{enumerate}
\item[(a)] Let $\rho \in (0, \infty]$ denote the injectivity radius of $M$, and define $T^\rho M$ to be the subset of $TM$ consisting of vectors of length less than $\rho$, and $D^\rho$ to be the subset $\{(p, q) : d_g(p, q) < \rho\} \subseteq M \times M$. Define $E : T^\rho M \to D^\rho$ by $E(x, v) = (x, \exp_x v)$. Prove that $E$ is a diffeomorphism.
\item[(b)] Use part (a) to prove that if $B$ is a topological space and $F, G : B \to M$ are continuous maps such that $d_g(F(x), G(x)) < \operatorname{inj}(M)$ for all $x \in B$, then $F$ and $G$ are homotopic.
\end{enumerate}
\end{exercise}

\begin{exercise}
\label{prob:k-point-homogeneous}
\dcref{ch6:6-18}
\uses{prob:inj-radius-diffeomorphism,prob:geodesic-completeness}

A connected Riemannian manifold $(M,g)$ is said to be $k$-point \emph{homogeneous} if for any two ordered $k$-tuples $(p_1,\ldots,p_k)$ and $(q_1,\ldots,q_k)$ of points in $M$ such that $d_g(p_i, p_j) = d_g(q_i,q_j)$ for all $i,j$, there is an isometry $\varphi: M \to M$ such that $\varphi(p_i) = q_i$ for $i = 1,\ldots,k$. Show that $(M,g)$ is 2-point homogeneous if and only if it is isotropic. [Hint: Assuming that $M$ is isotropic, first show that it is homogeneous by considering the midpoint of a geodesic segment joining sufficiently nearby points $p, q \in M$, and then use the result of Problem 6-12(c) to show that it is complete.] (Used on pp. 56, 261.)
\end{exercise}

\begin{exercise}
\label{prob:killing-vector-field-properties}
\dcref{ch6:6-24}
\uses{prob:killing-vector-field}

Let $(M, g)$ be a Riemannian manifold. Recall that a vector field $X \in \mathfrak{X}(M)$ is called a \emph{Killing vector field} if $\mathcal{L}_X g = 0$ (see Problem 5-22).
\begin{enumerate}
\item[(a)] Prove that a Killing vector field that is normal to a geodesic at one point is normal everywhere along the geodesic.
\item[(b)] Prove that if a Killing vector field vanishes at a point $p$, then it is tangent to geodesic spheres centered at $p$.
\item[(c)] Prove that a Killing vector field on an odd-dimensional manifold cannot have an isolated zero.
\end{enumerate}
(Used on p.\ 315.)
\end{exercise}

\begin{exercise}
\label{prob:lie-group-exp-surjective}
\dcref{ch6:6-13}
\uses{prob:lie-group-bi-invariant-metric}

Let $G$ be a connected compact Lie group. Show that the Lie group exponential map of $G$ is surjective. [Hint: Use Problem 5-8.]
\end{exercise}

\begin{exercise}
\label{prob:limit-isometry}
\dcref{ch6:6-8}

Suppose $(M, g)$ and $(\tilde{M}, \tilde{g})$ are connected Riemannian manifolds (not necessarily complete), and for each $i \in \mathbb{Z}^+$, $\varphi_i : M \to \tilde{M}$ is a Riemannian isometry such that $\varphi_i$ converges pointwise to a map $\varphi : M \to \tilde{M}$. Show that $\varphi$ is a Riemannian isometry. [Hint: Once you have shown that $\varphi$ is a local isometry, to show that $\varphi$ is surjective, suppose $y$ is a limit point of $\varphi(M)$. Choose $x \in M$ such that $\varphi(x)$ lies in a uniformly normal neighborhood of $y$, and show that there exists a convergent sequence of points $(x_i, v_i) \in TM$ such that $\varphi_i(x_i) = \varphi(x)$ and $\varphi_i(\exp_{x_i} v_i) = y$.]
\end{exercise}

\begin{exercise}
\label{prob:manifold-boundary-complete}
\dcref{ch6:6-22}

Let $(M,g)$ be a connected Riemannian manifold with boundary. Prove that $(M, d_g)$ is a complete metric space if and only if the following condition holds: for every geodesic $\gamma: [0,b) \to M$ that cannot be extended to a geodesic on any interval $[0,b')$ with $b' > b$, $\gamma(t)$ converges to a point of $\partial M$ as $t \nearrow b$.
\end{exercise}

\begin{exercise}[Characterization of metric isometry]
\label{prob:metric-isometry-characterization}
\label{prob:exponential-distance-limit}
\dcref{ch6:6-7}
\uses{prop:isometry-invariance-riemannian-distance,prob:distance-preserving-map-linear-isometry}

We now have two kinds of ``metrics'' on a connected Riemannian manifold: the Riemannian metric and the distance function. Correspondingly, there are two definitions of ``isometry'' between connected Riemannian manifolds: a \emph{Riemannian isometry} is a diffeomorphism that pulls one Riemannian metric back to the other, and a \emph{metric isometry} is a homeomorphism that preserves distances. Proposition 2.51 shows that every Riemannian isometry is a metric isometry. This problem outlines a proof of the converse. Suppose $(M, g)$ and $(\tilde{M}, \tilde{g})$ are connected Riemannian manifolds, and $\varphi: M \to \tilde{M}$ is a metric isometry.
\begin{enumerate}
\item[(a)] Show that for every $p \in M$ and $v, w \in T_p M$, we have
\[
\lim_{t \to 0} \frac{d_g(\exp_p tv, \exp_p tw)}{t} = |v - w|_g.
\]
[Hint: Use the Taylor series for $g$ in Riemannian normal coordinates on a convex geodesic ball centered at $p$.]
\item[(b)] Show that $\varphi$ takes geodesics to geodesics.
\item[(c)] For each $p \in M$, show that there exist an open ball $B_\epsilon(0) \subseteq T_p M$ and a continuous map $\psi : B_\epsilon(0) \to T_{\varphi(p)} \tilde{M}$ satisfying $\psi(0) = 0$ and $\exp_{\varphi(p)} \psi(v) = \varphi(\exp_p v)$ for all $v \in B_\epsilon(0)$.
\item[(d)] With $\psi$ as above, show that $|\psi(v) - \psi(w)|_{\tilde{g}} = |v - w|_g$ for all $v, w \in B_\epsilon(0)$, and conclude from Problem 2-2 that $\psi$ is the restriction of a linear isometry.
\item[(e)] With $p$ and $\psi$ as above, show that $\varphi$ is smooth on a neighborhood of $p$ and $d\varphi_p = \psi$.
\item[(f)] Conclude that $\varphi$ is a Riemannian isometry.
\end{enumerate}
\end{exercise}

\begin{exercise}
\label{prob:onion-skin-geodesics}
\dcref{ch6:6-15}

Let $S$ be the unit 2-sphere minus its north and south poles, and let $M = (0, \pi) \times \mathbb{R}$. Define $q : M \to S$ by $q(\varphi, \theta) = (\sin \varphi \cos \theta, \sin \varphi \sin \theta, \cos \varphi)$, and let $g$ be the metric on $M$ given by pulling back the round metric: $g = q^* \bar{g}$. (Think of $M$ as an infinitely long onion skin wrapping infinitely many times around an onion.) Prove that the Riemannian manifold $(M, g)$ has diameter $\pi$, but contains infinitely long properly embedded geodesics.
\end{exercise}

\begin{exercise}
\label{prob:parallel-vector-field}
\dcref{ch6:6-29}

Suppose $(M, g)$ is a Riemannian $n$-manifold that admits a nonzero parallel vector field $X$. Show that for each $p \in M$, there exist a neighborhood $U$ of $p$, a Riemannian $(n - 1)$-manifold $N$, and an open interval $I \subseteq \mathbb{R}$ with its Euclidean metric such that $U$ is isometric to the Riemannian product $N \times I$. [Hint: It might be helpful to prove that the 1-form $X^\flat$ is closed and to compute the Lie derivative $\mathcal{L}_X g$.]
\end{exercise}

\begin{exercise}
\label{prob:poincare-ball-distance-formula}
\dcref{ch6:6-3}
\uses{prob:sphere-hyperbolic-scaling,prob:poincare-disk}

Let $n$ be a positive integer and $R$ a positive real number. Prove that the Riemannian distance between any two points in the Poincaré ball model $(\mathbb{B}^n(R), \bar{g}_R)$ of hyperbolic space of radius $R$ is given by
\[
d_{\bar{g}_R}(p,q) = R \operatorname{arcosh}\left(1 + \frac{2R^2|p-q|^2}{(R^2 - |p|^2)(R^2 - |q|^2)}\right),
\]
where $|\cdot|$ represents the Euclidean norm in $\mathbb{R}^n$. [Hint: First use the result of Problem 3-5 to show that it suffices to consider the case $R = 1$. Then use a rotation to reduce to the case $n = 2$, and use the group action of Problem 3-8 to show that it suffices to consider the case in which $p$ is the origin.]
\end{exercise}

\begin{exercise}
\label{prob:ray-definition}
\dcref{ch6:6-21}

Let $(M,g)$ be a complete, connected, noncompact Riemannian manifold. Define a \emph{ray} in $M$ to be a geodesic $\gamma$ whose domain is $[0,\infty)$, and such that the restriction of $\gamma$ to $[0,b]$ is minimizing for every $b > 0$. Prove that for each $p \in M$ there is a ray in $M$ starting at $p$.
\end{exercise}

\begin{exercise}
\label{prob:smooth-function-equals-distance}
\dcref{ch6:6-27}

Suppose $(M, g)$ is a complete, connected Riemannian manifold, $S \subseteq M$ is a closed subset, and $f : M \to \mathbb{R}$ is a continuous function such that $f^{-1}(0) = S$ and $f$ is smooth with unit gradient on $M \setminus S$. Prove that $f(x) = d_g(x, S)$ for all $x \in M$.
\end{exercise}

\begin{exercise}
\label{prob:submanifold-distance-inequality}
\dcref{ch6:6-11}

Suppose $(M, g)$ is a connected Riemannian manifold, $P \subseteq M$ is a connected embedded submanifold, and $\bar{g}$ is the induced Riemannian metric on $P$.
\begin{enumerate}
\item[(a)] Show that $d_{\bar{g}}(p, q) \geq d_g(p, q)$ for $p, q \in P$.
\item[(b)] Prove that if $(M, g)$ is complete and $P$ is closed in $M$, then $(P, \bar{g})$ is complete.
\item[(c)] Give an example of a complete Riemannian manifold $(M, g)$ and a connected embedded submanifold $P \subseteq M$ that is complete but not closed in $M$.
\end{enumerate}
\end{exercise}

\begin{exercise}
\label{prob:symmetric-space-homogeneous}
\dcref{ch6:6-19}
\uses{prob:k-point-homogeneous}

Prove that every Riemannian symmetric space is homogeneous. [Hint: Proceed as in Problem 6-18.] (Used on p. 78.)
\end{exercise}

\begin{exercise}
\label{prob:tubular-neighborhood-volume}
\dcref{ch6:6-31}
\uses{prob:volume-geodesic-ball}

Suppose $(M, g)$ is a Riemannian $n$-manifold, $P \subseteq M$ is a compact embedded submanifold, and $U$ is an $\varepsilon$-tubular neighborhood of $P$ for some $\varepsilon > 0$. For $0 < \delta < \varepsilon$, define $U_\delta$ and $P_\delta$ by
\[
U_\delta = \{x \in U : d_g(x, P) \leq \delta\},
\]
\[
P_\delta = \{x \in U : d_g(x, P) = \delta\}.
\]
\begin{enumerate}
\item[(a)] Prove that $U_\delta$ is a regular domain and $P_\delta$ is a compact, embedded submanifold of $M$ for each $\delta \in (0, \varepsilon)$.
\item[(b)] Generalize the result of Problem 6-30 by proving that
\[
\operatorname{Vol}(U_\delta) = \int_0^\delta \operatorname{Area}(P_r) \, dr,
\]
where $\operatorname{Area}(P_r)$ denotes the $(n-1)$-dimensional volume of $P_r$ with the induced Riemannian metric.
\end{enumerate}
\end{exercise}

\begin{exercise}
\label{prob:volume-geodesic-ball}
\dcref{ch6:6-30}

Suppose $(M, g)$ is a Riemannian $n$-manifold, $p \in M$, and $B_\varepsilon(p)$ is a geodesic ball centered at $p$. Prove that for every $\delta$ such that $0 < \delta < \varepsilon$,
\[
\operatorname{Vol}\left(\overline{B_\delta(p)}\right) = \int_0^\delta \operatorname{Area}(\partial B_r(p)) \, dr,
\]
where $\operatorname{Area}(\partial B_r(p))$ represents the $(n-1)$-dimensional volume of $\partial B_r(p)$ with its induced Riemannian metric.
\end{exercise}
